\documentclass[11pt,a4paper]{article}
\input{../../sablony/pripravy/priprava.sty}
\usepackage{tikz,pgfplots}
\pgfplotsset{compat=1.18}
\nastavPripravu{1. ročník}{Kvadratické rovnice a nerovnice}{120}{Slovní úlohy na kvadratické rovnice II}
\begin{document}
\blok{Slovní úlohy na kvadratické rovnice II}
Technické aplikace: brzdná dráha, trajektorie, výrobní optimalizace s modelovými parametry.
\blok{Metoda a návaznost na graf}
\begin{itemize}
\item V každé úloze uveď jednotky, zjednodušující předpoklady a konečný smysluplný výsledek.
\item Číselné provozní parametry jsou ilustrativní modelová data, nikoli deklarovaná měření.
\end{itemize}
\blok{Společně řešené příklady}
\begin{prikladReseniBox}[title={Příklad 1}]
Dráha brzdění je v modelu $s=\frac{v^2}{2a}$, $a=5\,\mathrm{m/s^2}$. Jakou počáteční rychlost (v m/s) odpovídá dráha $40\,\mathrm m$?
\tcblower
\[
\begin{aligned}
40&=\frac{v^2}{10}\\[2pt]
v^2&=400\\[2pt]
v&=20\,\mathrm{m/s}\quad(v\ge0)
\end{aligned}
\]
\end{prikladReseniBox}
\begin{prikladReseniBox}[title={Příklad 2}]
Výnos výrobního procesu je modelován $Z(q)=-2q^2+40q-150$ (tis. Kč). Pro které nezáporné $q$ je $Z(q)>0$?
\tcblower
\[
\begin{aligned}
-2q^2+40q-150&>0\\[2pt]
q^2-20q+75&<0\\[2pt]
(q-5)(q-15)&<0\\[2pt]
K&=(5;15)
\end{aligned}
\]
\end{prikladReseniBox}
\blok{Samostatné procvičování}
\begin{zadaniReseniBox}[title={Samostatná práce 1}]
Obdélníková obrazovka má obsah $96\,\mathrm{cm^2}$, šířka je o $4\,\mathrm{cm}$ menší než výška. Určete rozměry.
\tcblower
\[
\begin{aligned}
x(x-4)&=96\\[2pt]
x^2-4x-96&=0\\[2pt]
(x-12)(x+8)&=0\\[2pt]
\text{Výška}&=12\,\mathrm{cm},\ \text{šířka}=8\,\mathrm{cm}
\end{aligned}
\]
\end{zadaniReseniBox}
\begin{zadaniReseniBox}[title={Samostatná práce 2}]
Model výkonové ztráty odporu je $P=I^2R$, $R=8\,\Omega$. Pro jaký nezáporný proud dosahuje $P=72\,\mathrm W$?
\tcblower
\[
\begin{aligned}
72&=8I^2\\[2pt]
I^2&=9\\[2pt]
I&=3\,\mathrm A\quad(I\ge0)
\end{aligned}
\]
\end{zadaniReseniBox}
\blok{Volitelná rezerva}
\poznamka{Studenti si vymyslí obdobný příklad se změnou jednoho koeficientu a vysvětlí, jak se změní počet kořenů či řešení.}
\blok{Rozšiřující příklady ze sbírky}
\poznamka{Využij jako kontrast ke kvadratickým rovnicím: hledání extrému může být rychlejší ve vrcholovém tvaru.}

\begin{zadaniReseniBox}[title={Rozšíření 1}]
Určete maximum/minimum a bod extrému výrazu: $x \cdot \left(1 - 2x\right)$.
\tcblower
\[
\begin{aligned}x_v&=-\frac{1}{2\cdot(-2)}=\frac{1}{4}\\\max f&=f(\frac{1}{4})=\frac{1}{8}\\V&=[\frac{1}{4};\frac{1}{8}]\end{aligned}
\]
\end{zadaniReseniBox}

\begin{zadaniReseniBox}[title={Rozšíření 2}]
Určete maximum/minimum a bod extrému výrazu: $- x^{2} + 2x + 3$.
\tcblower
\[
\begin{aligned}x_v&=-\frac{2}{2\cdot(-1)}=1\\\max f&=f(1)=4\\V&=[1;4]\end{aligned}
\]
\end{zadaniReseniBox}

\blok{Volitelná výzva -- vrchol místo derivace}
\begin{zadaniReseniBox}[title={Rozšíření 3}]
Určete maximum $p(x)=-(x-2)^2+9$ a hodnotu $x$, kdy nastává.
\tcblower
\[\begin{aligned}
p(x)&=9-(x-2)^2\le9\\x&=2\\\max p&=9
\end{aligned}\]
\end{zadaniReseniBox}
\blok{Domácí úkol}
\begin{zadaniReseniBox}[title={Domácí úkol}]
\begin{enumerate}[label=\alph*),itemsep=2pt]
\item Obdélník má obsah $60\,\mathrm{cm^2}$ a delší strana je o $7\,\mathrm{cm}$ větší. Určete délky stran.
\item Čtvercový pozemek má po rozšíření každé strany o $3\,\mathrm m$ obsah $100\,\mathrm{m^2}$. Jaká byla původní strana?
\item Model výšky míčku je $h(t)=-5t^2+20t+1$ (metry, sekundy). Kdy je míček ve výšce $16\,\mathrm m$?
\item Dráha brzdění je v modelu $s=\frac{v^2}{2a}$, $a=5\,\mathrm{m/s^2}$. Jakou počáteční rychlost (v m/s) odpovídá dráha $40\,\mathrm m$?
\item Výnos výrobního procesu je modelován $Z(q)=-2q^2+40q-150$ (tis. Kč). Pro které nezáporné $q$ je $Z(q)>0$?
\item Obdélníková obrazovka má obsah $96\,\mathrm{cm^2}$, šířka je o $4\,\mathrm{cm}$ menší než výška. Určete rozměry.
\end{enumerate}
\tcblower
\begin{enumerate}[label=\alph*),itemsep=2pt]
\item
\[\begin{aligned}
x(x+7)&=60\\[2pt]
x^2+7x-60&=0\\[2pt]
(x-5)(x+12)&=0\\[2pt]
x&=5\,\mathrm{cm}\quad (x=-12\text{ nelze})\\[2pt]
\text{Strany: }&5\,\mathrm{cm};\ 12\,\mathrm{cm}
\end{aligned}\]
\item
\[\begin{aligned}
(x+3)^2&=100\\[2pt]
x+3&=10\quad (x>0)\\[2pt]
x&=7\,\mathrm m
\end{aligned}\]
\item
\[\begin{aligned}
-5t^2+20t+1&=16\\[2pt]
t^2-4t+3&=0\\[2pt]
(t-1)(t-3)&=0\\[2pt]
t&=1\,\mathrm s\ \text{nebo}\ 3\,\mathrm s
\end{aligned}\]
\item
\[\begin{aligned}
40&=\frac{v^2}{10}\\[2pt]
v^2&=400\\[2pt]
v&=20\,\mathrm{m/s}\quad(v\ge0)
\end{aligned}\]
\item
\[\begin{aligned}
-2q^2+40q-150&>0\\[2pt]
q^2-20q+75&<0\\[2pt]
(q-5)(q-15)&<0\\[2pt]
K&=(5;15)
\end{aligned}\]
\item
\[\begin{aligned}
x(x-4)&=96\\[2pt]
x^2-4x-96&=0\\[2pt]
(x-12)(x+8)&=0\\[2pt]
\text{Výška}&=12\,\mathrm{cm},\ \text{šířka}=8\,\mathrm{cm}
\end{aligned}\]
\end{enumerate}
\end{zadaniReseniBox}
\poznamka{Orientační čas: úvod 5 min, společné příklady 18 min, samostatně 15 min, kontrola a shrnutí 7 min. Rezerva je dobrovolná. Číselné parametry modelových aplikací nejsou skutečně naměřená data.}
\end{document}
